\documentclass[10pt,a4paper]{amsart}
\usepackage[margin=19mm]{geometry}
\usepackage{amsmath,amssymb,mathtools}
\usepackage[scr=rsfs,cal=euler]{mathalfa}
\usepackage{xfrac}
\usepackage[unicode,hidelinks,bookmarks=false]{hyperref}
\newcommand{\ie}{i.e.\ }
\newcommand{\cf}{cf.\ }
\newcommand{\resp}{respectively\ }
\newcommand{\Subsection}{\subsection}
\providecommand{\nobreakdash}{\nobreak\hbox{-}}
\setlength{\emergencystretch}{2em}
\multlinegap0pt
\usepackage{bm}
\usepackage{bbm}
\usepackage{dsfont}
\usepackage{xspace}
\usepackage{cancel}
\usepackage{mathtools}
\usepackage{amsthm}
\makeatletter
\@ifundefined{theorem}{\newtheorem{theorem}{Theorem}}{}
\@ifundefined{lemma}{\newtheorem{lemma}[theorem]{Lemma}}{}
\@ifundefined{proposition}{\newtheorem{proposition}[theorem]{Proposition}}{}
\makeatother
\newcommand{\abs}[1]{\left|{#1}\right|}
\newcommand{\bb}{\mathbb}
\newcommand{\bff}{\boldsymbol}
\newcommand{\curlh}{\mathrm{curl}_h\,}
\newcommand{\curl}{\mathrm{curl}\,}
\newcommand{\divg}{\mathrm{div}\,}
\newcommand{\inpro}[2]{\left\langle#1,#2\right\rangle}
\newcommand{\norm}[2]{\left\|{#1}\right\|_{#2}}
\title[Stochastic resistive Hall--MHD with current fluctuations: ma]{Stochastic resistive Hall--MHD with current fluctuations: martingale weak solutions via a convergent structure-preserving finite element method}
\author{Agus L. Soenjaya}
\date{Source: arXiv:2608.13407v1 (CC BY 4.0)}
\begin{document}
\maketitle

\noindent\emph{Source and licence.} Stochastic resistive Hall--MHD with current fluctuations: martingale weak solutions via a convergent structure-preserving finite element method. Version arXiv:2608.13407v1 (CC BY 4.0), distributed under the Creative Commons Attribution 4.0 International licence, as recorded in the arXiv licence metadata for this exact version and shown on the abstract page https://arxiv.org/abs/2608.13407v1. Full rights notice: licenses/arxiv-2608.13407-CC-BY-4.0.txt.\par\smallskip

\noindent\textit{Editorial scope.} Counted unit: the Proposition whose source label is \texttt{prop:stoch stability revised} in the article (source0.tex lines 1072--1143, source label \texttt{prop:stoch stability revised}), delivered together with the article's own complete original proof of it (source0.tex lines 1145--1706, one \texttt{proof} environment of 562 lines). No mathematics is written, repaired or re-derived: statement and proof are the article's own text line for line. Required context retained uncounted: equ:chi assumption (lines 445--448); equ:g growth (lines 452--463); equ:hall noise small (lines 465--467); equ:stoch fem (lines 1005--1040); equ:stoch fem B (lines 1006--1039); equ:stoch fem Ohm (lines 1006--1039); equ:hall current coeff bound final (lines 1058--1062); lem:weighted-discrete-holder (lines 4427--4438); eq:weighted-discrete-holder (lines 4429--4436); lem:discrete-BDG (lines 4441--4466). Every definition, display and label of those lines is retained unchanged. Omitted from this package and not redistributed: 1--444, 449--451, 464--464, 468--1004, 1040--1057, 1063--1071, 1144--1144, 1707--4426, 4437--4440, 4467--5547. Nothing inside a retained excerpt is omitted or altered: every retained body is delivered exactly as the article's lines are written, so no omission receipt of any kind accompanies this package. Citations: the retained excerpts carry no citation command at all, so no bibliography entry of the article is transcribed and no reference is redistributed. Reference adaptation: none was needed and none was made; every reference inside the delivered unit resolves inside it, so no unresolved reference or citation is produced. Printed slips kept literal: no printed word, symbol, hypothesis or number of the counted statement or of its proof was altered. Layout and class: the article's own class is \texttt{amsart} with packages including amsmath, amsthm, amscd, amsfonts, amssymb, graphicx, color, mathtools, mathrsfs, hyperref, enumerate, setspace, multicol, geometry; the wrapper is the lane's pinned \texttt{amsart} editorial wrapper at 10pt on A4, which restates the theorem-like environments the retained text needs and adds the article's own macro definitions that the retained text needs (\texttt{\textbackslash abs}, \texttt{\textbackslash bb}, \texttt{\textbackslash bff}, \texttt{\textbackslash curl}, \texttt{\textbackslash curlh}, \texttt{\textbackslash divg}, \texttt{\textbackslash inpro}, \texttt{\textbackslash norm}), with extra wrapper packages (bm, bbm, dsfont, xspace, cancel, mathtools). The delivered numbering is the unit's own. Assets: no retained span contains a figure environment, a graphics-inclusion command, a PDF-page inclusion, a table or a TikZ picture; no image file is copied into this package, the package declares no asset dependency and no image file is redistributed with it. Licence: version arXiv:2608.13407v1 of this article is distributed under the Creative Commons Attribution 4.0 International licence, as recorded in the arXiv licence metadata for this exact version and shown on the abstract page https://arxiv.org/abs/2608.13407v1. A full rights notice is delivered at licenses/arxiv-2608.13407-CC-BY-4.0.txt. Characters: every retained excerpt is pure ASCII, so the delivered engine, XeLaTeX with its default Latin Modern text font, reports no missing character.
\begin{align}\label{equ:chi assumption}
	\bff{\chi}_i\in \bb{W}^{1,\infty}_0(\mathscr{D}),\qquad
	\sum_{i=1}^\infty \norm{\bff{\chi}_i}{\bb{W}^{1,\infty}}^2<\infty.
\end{align}

\begin{align}
	\label{equ:g growth}
	\sum_{i=1}^\infty \norm{\bff{g}_i(\bff{v},\bff{C})}{\bb{L}^2}^2
	&\leq
	C\bigl(1+\norm{\bff{v}}{\bb{L}^2}^2+\norm{\bff{C}}{\bb{L}^2}^2\bigr), \quad \forall \bff{v},\bff{C}\in \bb{L}^2,
	\\
	\label{equ:g lipschitz}
	\sum_{i=1}^\infty \norm{\bff{g}_i(\bff{v}_1,\bff{C}_1)-\bff{g}_i(\bff{v}_2,\bff{C}_2)}{\bb{L}^2}^2
	&\leq
	L\bigl(\norm{\bff{v}_1-\bff{v}_2}{\bb{L}^2}^2+\norm{\bff{C}_1-\bff{C}_2}{\bb{L}^2}^2\bigr),
	\quad \forall \bff{v}_1,\bff{v}_2,\bff{C}_1,\bff{C}_2\in \bb{L}^2.
\end{align}

\begin{align}\label{equ:hall noise small}
	\delta_\chi:=C_{\mathscr{D}}\eta^2\sum_{i=1}^\infty \norm{\bff{\chi}_i}{\bb{W}^{1,\infty}}^2<\sigma,
\end{align}

\begin{subequations}\label{equ:stoch fem}
	\begin{align}
		&\inpro{\bff{u}_h^n-\bff{u}_h^{n-1}}{\bff{\phi}_h}
		+\nu\tau\inpro{\nabla\bff{u}_h^n}{\nabla\bff{\phi}_h}
		+\frac{\tau}{2}\Big[\inpro{(\bff{u}_h^{n-1}\cdot\nabla)\bff{u}_h^n}{\bff{\phi}_h}-\inpro{(\bff{u}_h^{n-1}\cdot\nabla)\bff{\phi}_h}{\bff{u}_h^n}\Big]
		\nonumber\\
		&\quad
		-\tau\inpro{p_h^n}{\divg\bff{\phi}_h}
		-\tau\inpro{\bff{J}_h^n\times\bff{B}_h^{n-1}}{\bff{\phi}_h}
		=
		\tau\inpro{\bff{f}^n}{\bff{\phi}_h}
		+\sum_{i=1}^\infty\inpro{\bff{g}_i(\bff{u}_h^{n-1},\bff{B}_h^{n-1})}{\bff{\phi}_h}\Delta_n\beta_i^u,
		\label{equ:stoch fem u}
		\\[1ex]
		&\inpro{\divg\bff{u}_h^n}{q_h}=0,
		\label{equ:stoch fem div u}
		\\[1ex]
		&\inpro{\bff{B}_h^n-\bff{B}_h^{n-1}}{\bff{\psi}_h}
		+\tau\inpro{\curl\bff{E}_h^n}{\bff{\psi}_h}
		=
		-\sum_{i=1}^{\infty}
		\inpro{\curl \Pi_h^{\bb{X}}\bigl(\sigma\bff{\chi}_i+\eta\bff{\chi}_i\times\mathcal{R}_h\bff{B}_h^{n-1}\bigr)}{\bff{\psi}_h}\Delta_n\beta_i^B.
		\label{equ:stoch fem B}
		\\[1ex]
		&\sigma\inpro{\bff{J}_h^n}{\bff{\chi}_h}
		+\eta\inpro{\bff{J}_h^n\times\bff{B}_h^{n-1}}{\bff{\chi}_h}
		=
		\inpro{\bff{E}_h^n}{\bff{\chi}_h}
		+\inpro{\bff{u}_h^n\times\bff{B}_h^{n-1}}{\bff{\chi}_h},
		\label{equ:stoch fem Ohm}
		\\[1ex]
		&\inpro{\bff{J}_h^n}{\bff{\omega}_h}
		-\inpro{\bff{B}_h^n}{\curl\bff{\omega}_h}=0,
		\label{equ:stoch fem J}
	\end{align}
\end{subequations}

\begin{align}\label{equ:hall current coeff bound final}
	\sum_{i=1}^{\infty}\norm{\curl \Pi_h^{\bb{X}}\bigl(\sigma\bff{\chi}_i+\eta\bff{\chi}_i\times\mathcal{R}_h\bff{B}_h^{n-1}\bigr)}{\bb{L}^2}^2
	\leq
	C_\chi\bigl(1+\norm{\bff{B}_h^{n-1}}{\bb{L}^2}^2\bigr)+\delta_\chi\norm{\bff{J}_h^{n-1}}{\bb{L}^2}^2,
\end{align}

\begin{lemma}[The discrete Jensen inequality]\label{lem:weighted-discrete-holder}
	Let \(q\ge1\) and $k>0$. let \(\{a_j\}_{j=0}^{J-1}\) be a non-negative sequence and set $T=Jk$. Then
	\begin{equation}\label{eq:weighted-discrete-holder}
		\left(
		\sum_{j=0}^{J-1} k a_j
		\right)^q
		\le
		T^{q-1}
		\sum_{j=0}^{J-1} k a_j^q.
	\end{equation}
	The sign in \eqref{eq:weighted-discrete-holder} is reversed if $q\in (0,1)$.
\end{lemma}

\begin{lemma}[The discrete Burkholder--Davis--Gundy (BDG) inequality]\label{lem:discrete-BDG}
	Let $p\ge 1$, let $H$ be a Hilbert space, and let
	$\{\bff Z_j\}_{j=0}^{J-1}$ be an $H$-valued martingale difference sequence
	with respect to a filtration \(\{\mathcal F_j\}_{j=0}^{J}\), that is,
	\(\bff Z_j\) is \(\mathcal F_{j+1}\)-measurable and
	\[
	\mathbb E[\bff Z_j\mid\mathcal F_j]=\bff0,
	\qquad
	j=0,\ldots,J-1.
	\]
	Assume that \(\bff Z_j\in L^p(\Omega;H)\) for each \(j\). Then there exists
	a constant \(C_p>0\), depending only on \(p\), such that
	\begin{equation}\label{eq:discrete-BDG-square-function}
		\bb E\left[
		\max_{0\le r\le J}
		\norm{\sum_{j=0}^{r-1}\bff Z_j}{H}^{p}
		\right]
		\le
		C_p\,
		\bb E\left[
		\left(
		\sum_{j=0}^{J-1}\norm{\bff Z_j}{H}^{2}
		\right)^{p/2}
		\right].
	\end{equation}
\end{lemma}

\begin{proposition}
	\label{prop:stoch stability revised}
	Let $(\bff{u}_h^n,p_h^n,\bff{B}_h^n,\bff{E}_h^n,\bff{J}_h^n)$, $n=0,\ldots,N$, be a solution of the fully discrete scheme \eqref{equ:stoch fem}. Suppose that
	\[
	\divg\bff{B}_h^0=0,
	\qquad
	\Pi_h^{\bb{X}} \bff{J}_h^0=\curlh\bff{B}_h^0,
	\]
	and that the velocity-noise growth condition \eqref{equ:g growth} and the current-noise assumptions \eqref{equ:chi assumption} and \eqref{equ:hall noise small} hold.
	Then
	\begin{align}\label{equ:div B exact revised}
		\divg\bff{B}_h^n=0
		\quad\text{pointwise on each element of }\mathcal T_h,
		\qquad n=0,\ldots,N.
	\end{align}
	Moreover, for every integer $p\ge1$, there exists
	$\delta_{p}^\ast>0$ such that, if
	$\delta_\chi\le \delta_{p}^\ast$, then
	\begin{align}\label{equ:stoch stability revised main}
		&\bb E \left[\max_{0\le n\le N}
		\left(
		\norm{\bff u_h^n}{\bb L^2}^2
		+
		\norm{\bff B_h^n}{\bb L^2}^2
		\right)^p\right]
		+
		\bb E\left[
		\left(
		\sum_{n=1}^{N}
		\norm{\bff u_h^n-\bff u_h^{n-1}}{\bb L^2}^2
		+
		\sum_{n=1}^{N}
		\norm{\bff B_h^n-\bff B_h^{n-1}}{\bb L^2}^2
		\right)^p
		\right]
		\nonumber\\
		&\quad
		+
		\bb E\left[
		\left(
		\sum_{n=1}^{N}\tau
		\left(
		\nu\norm{\nabla\bff u_h^n}{\bb L^2}^2
		+
		\sigma\norm{\bff J_h^n}{\bb L^2}^2
		\right)
		\right)^p
		\right]
		\nonumber\\
		&\le
		C_{p,T}
		\left(
		1
		+
		\bb E\left[
		\left(
		\norm{\bff u_h^0}{\bb L^2}^2
		+
		\norm{\bff B_h^0}{\bb L^2}^2
		\right)^p
		\right]
		+
		\bb E\left[
		\left(
		\tau\norm{\bff J_h^0}{\bb L^2}^2
		\right)^p
		\right]
		+
		\norm{\bff f}{L^2_T(\bb L^2)}^{2p}
		\right).
	\end{align}
\end{proposition}

\begin{proof}
	We divide the proof into several steps.
	
	\emph{Step 1: Exact preservation of magnetic divergence.}
	Since
	\[
	\bff E_h^n\in\bb X_h^0,
	\qquad
	\Pi_h^{\bb X}\bigl(\sigma\bff\chi_i+\eta\bff\chi_i\times\mathcal R_h\bff B_h^{n-1}\bigr)\in\bb X_h^0,
	\]
	we have
	\[
	\curl\bff E_h^n\in\bb{RT}_h^0,
	\qquad
	\curl\Pi_h^{\bb X}\bigl(\sigma\bff\chi_i+\eta\bff\chi_i\times\mathcal R_h\bff B_h^{n-1}\bigr)\in\bb{RT}_h^0.
	\]
	The magnetic update \eqref{equ:stoch fem B} therefore gives, as an identity in $\bb{RT}_h^0$,
	\begin{align}\label{equ:B update strong revised}
		\bff B_h^n
		=
		\bff B_h^{n-1}
		-\tau\,\curl\bff E_h^n
		-\sum_{i=1}^{\infty}
		\curl\Pi_h^{\bb X}\bigl(\sigma\bff\chi_i+\eta\bff\chi_i\times\mathcal R_h\bff B_h^{n-1}\bigr)\Delta_n\beta_i^B.
	\end{align}
	Taking the elementwise divergence and using $\divg\curl\bff\omega_h=0$ for all $\bff\omega_h\in\bb X_h^0$, we get
	\[
	\divg\bff B_h^n=\divg\bff B_h^{n-1}.
	\]
	The claim \eqref{equ:div B exact revised} follows by induction over $n$.
	
	\emph{Step 2: Discrete energy identity.}
	We define $\bff g_{i,h}^{n-1}:= \Pi_h^{\bb V} \bff g_i(\bff u_h^{n-1},\bff B_h^{n-1})$ and set
	\begin{align}\label{equ:G u G B revised}
		\bff G_{u,h}^n
		:=
		\sum_{i=1}^{\infty}\bff g_{i,h}^{n-1}\Delta_n\beta_i^u,
		\qquad
		\bff G_{B,h}^n
		:=
		\sum_{i=1}^{\infty}
		\Pi_h^{\bb X}\bigl(\sigma\bff\chi_i+\eta\bff\chi_i\times\mathcal R_h\bff B_h^{n-1}\bigr)\Delta_n\beta_i^B.
	\end{align}
	For ease of presentation, we define the discrete energy
	\[
	\mathcal E_h^n
	:=
	\frac12\norm{\bff u_h^n}{\bb L^2}^2
	+
	\frac12\norm{\bff B_h^n}{\bb L^2}^2,
	\]
	and take 
	\[
	\bff\phi_h=\bff u_h^n,\qquad
	q_h=p_h^n,\qquad
	\bff\psi_h=\bff B_h^n,\qquad
	\bff\chi_h=\bff J_h^n,\qquad
	\bff\omega_h=\bff E_h^n
	\]
	in the scheme \eqref{equ:stoch fem}. We note the identities
	\begin{align*}
	\inpro{\curl\bff E_h^n}{\bff B_h^n}
	&=
	\inpro{\bff E_h^n}{\bff J_h^n}
	=
	\sigma\norm{\bff J_h^n}{\bb L^2}^2
	-
	\inpro{\bff u_h^n\times\bff B_h^{n-1}}{\bff J_h^n},
	\end{align*}
	by the definition of $\curlh$ and Ohm's law \eqref{equ:stoch fem Ohm}. 
	Using the elementary vector identity
	\[
	\inpro{\bff a-\bff b}{\bff a}
	=
	\frac12\norm{\bff a}{\bb L^2}^2
	-
	\frac12\norm{\bff b}{\bb L^2}^2
	+
	\frac12\norm{\bff a-\bff b}{\bb L^2}^2,
	\]
	we obtain the pathwise energy identity
	\begin{align}\label{equ:energy identity revised}
		&\mathcal E_h^n-\mathcal E_h^{n-1}
		+
		\frac12\norm{\bff u_h^n-\bff u_h^{n-1}}{\bb L^2}^2
		+
		\frac12\norm{\bff B_h^n-\bff B_h^{n-1}}{\bb L^2}^2
		+
		\nu\tau\norm{\nabla\bff u_h^n}{\bb L^2}^2
		+
		\sigma\tau\norm{\bff J_h^n}{\bb L^2}^2
		\nonumber\\
		&\qquad
		=
		\tau\inpro{\bff f^n}{\bff u_h^n}
		+
		\inpro{\bff G_{u,h}^n}{\bff u_h^n}
		-
		\inpro{\curl\bff G_{B,h}^n}{\bff B_h^n}.
	\end{align}
	
	\emph{Step 3: Pathwise inequality with martingale increment.}
	Split
	\[
	\inpro{\bff G_{u,h}^n}{\bff u_h^n}
	=
	\inpro{\bff G_{u,h}^n}{\bff u_h^{n-1}}
	+
	\inpro{\bff G_{u,h}^n}{\bff u_h^n-\bff u_h^{n-1}},
	\]
	and
	\[
	-\inpro{\curl\bff G_{B,h}^n}{\bff B_h^n}
	=
	-\inpro{\curl\bff G_{B,h}^n}{\bff B_h^{n-1}}
	-
	\inpro{\curl\bff G_{B,h}^n}{\bff B_h^n-\bff B_h^{n-1}}.
	\]
	Define
	\begin{align}\label{equ:mart increment revised}
		M_h^n
		:=
		\inpro{\bff G_{u,h}^n}{\bff u_h^{n-1}}
		-
		\inpro{\curl\bff G_{B,h}^n}{\bff B_h^{n-1}}.
	\end{align}
	Young's inequality then yields
	\begin{align}\label{equ:pathwise inequality revised}
		&\mathcal E_h^n-\mathcal E_h^{n-1}
		+
		c_0\Big(
		\norm{\bff u_h^n-\bff u_h^{n-1}}{\bb L^2}^2
		+
		\norm{\bff B_h^n-\bff B_h^{n-1}}{\bb L^2}^2
		+
		\tau\norm{\nabla\bff u_h^n}{\bb L^2}^2
		+
		\tau\norm{\bff J_h^n}{\bb L^2}^2
		\Big)
		\nonumber\\
		&\qquad
		\le
		C\tau\norm{\bff f^n}{\widetilde{\bb H}^{-1}}^2
		+
		M_h^n
		+
		C\norm{\bff G_{u,h}^n}{\bb L^2}^2
		+
		C\norm{\curl\bff G_{B,h}^n}{\bb L^2}^2,
	\end{align}
	where $c_0>0$ depends on $\nu$ and $\sigma$ but not on $h,\tau,n$.
	
	\emph{Step 4: First moment estimate.}
	We note that the conditional It\^o isometry gives
	\begin{align}\label{equ:cond Ito u revised}
		\bb E\left[
		\norm{\bff G_{u,h}^n}{\bb L^2}^2\mid\mathscr F_{t_{n-1}}
		\right]
		=
		\tau\sum_{i=1}^{\infty}\norm{\bff g_{i,h}^{n-1}}{\bb L^2}^2,
	\end{align}
	and
	\begin{align}\label{equ:cond Ito B revised}
		\bb E\left[
		\norm{\curl\bff G_{B,h}^n}{\bb L^2}^2\mid\mathscr F_{t_{n-1}}
		\right]
		&=
		\tau\sum_{i=1}^{\infty}
		\norm{
			\curl\Pi_h^{\bb X}\bigl(\sigma\bff\chi_i+\eta\bff\chi_i\times\mathcal R_h\bff B_h^{n-1}\bigr)
		}{\bb L^2}^2
		\nonumber\\
		&\le
		C_\chi \tau \bigl(1+\norm{\bff{B}_h^{n-1}}{\bb{L}^2}^2\bigr)+\delta_\chi \tau \norm{\bff{J}_h^{n-1}}{\bb{L}^2}^2,
	\end{align}
	where we also used \eqref{equ:hall current coeff bound final}.
	
	
	Now, taking conditional expectation in \eqref{equ:pathwise inequality revised}, noting the fact that $\bb E[M_h^n\mid\mathscr F_{t_{n-1}}]=0$, using \eqref{equ:cond Ito u revised}--\eqref{equ:cond Ito B revised}, \eqref{equ:g growth}, and the tower property, we obtain
	\begin{align}\label{equ:first moment one step revised}
		&\bb E \left[\mathcal E_h^n\right]
		+
		c_0\bb E\Big[
		\norm{\bff u_h^n-\bff u_h^{n-1}}{\bb L^2}^2
		+
		\norm{\bff B_h^n-\bff B_h^{n-1}}{\bb L^2}^2
		+
		\tau\norm{\nabla\bff u_h^n}{\bb L^2}^2
		+
		\tau\norm{\bff J_h^n}{\bb L^2}^2
		\Big]
		\nonumber\\
		&\qquad
		\le
		(1+C\tau)\bb E\mathcal E_h^{n-1}
		+
		C\tau
		+
		C\tau\bb E\norm{\bff f^n}{\widetilde{\bb H}^{-1}}^2
		+
		C\delta_\chi\tau\bb E\norm{\bff J_h^{n-1}}{\bb L^2}^2.
	\end{align}
	Summing over $n=1,\ldots,m$, and absorbing the last term to the left-hand side when $\delta_\chi$ is sufficiently small, we obtain the first moment estimate by the discrete Gronwall lemma.
	
	
	\emph{Step 5: Quadratic stochastic remainders.}
	For $m=1,\ldots,N$, set
	\[
	Q_{u,m}
	:=
	\sum_{n=1}^{m}
	\norm{\bff G_{u,h}^n}{\bb L^2}^2,
	\qquad
	Q_{B,m}
	:=
	\sum_{n=1}^{m}
	\norm{\curl\bff G_{B,h}^n}{\bb L^2}^2.
	\]
	Conditionally on $\mathscr F_{t_{n-1}}$, the random variables
	$\bff G_{u,h}^n$ and $\curl\bff G_{B,h}^n$ are centred
	$\bb L^2$-valued Gaussian random variables. Hence, for every integer
	$p\ge1$,
	\begin{align}
		\label{equ:conditional Gaussian Gu}
		\bb E\left[
		\norm{\bff G_{u,h}^n}{\bb L^2}^{2p}
		\mid\mathscr F_{t_{n-1}}
		\right]
		&\le
		C_p\tau^p
		\left(
		\sum_{i=1}^{\infty}
		\norm{\bff g_{i,h}^{n-1}}{\bb L^2}^2
		\right)^p,
		\\
		\label{equ:conditional Gaussian GB}
		\bb E\left[
		\norm{\curl\bff G_{B,h}^n}{\bb L^2}^{2p}
		\mid\mathscr F_{t_{n-1}}
		\right]
		&\le
		C_p\tau^p
		\left(
		\sum_{i=1}^{\infty}
		\norm{
			\curl\Pi_h^{\bb X}
			\bigl(
			\sigma\bff\chi_i
			+
			\eta\bff\chi_i\times\mathcal R_h\bff B_h^{n-1}
			\bigr)}
		{\bb L^2}^2
		\right)^p.
	\end{align}
	The corresponding identities for $p=1$ are
	\eqref{equ:cond Ito u revised} and \eqref{equ:cond Ito B revised}.
	
	For $p>1$, the elementary inequality
	\[
	(x+y)^p-x^p
	\le
	C_px^{p-1}y+C_py^p,
	\qquad x,y\ge0,
	\]
	gives
	\[
	Q_{u,n}^p-Q_{u,n-1}^p
	\le
	C_pQ_{u,n-1}^{p-1}
	\norm{\bff G_{u,h}^n}{\bb L^2}^2
	+
	C_p\norm{\bff G_{u,h}^n}{\bb L^2}^{2p}.
	\]
	Taking conditional expectations, summing over $n=1,\ldots,m$, and using
	$Q_{u,n-1}\le Q_{u,m}$, we obtain
	\begin{align*}
		\bb E[Q_{u,m}^p]
		&\le
		C_p\bb E\left[
		Q_{u,m}^{p-1}
		\sum_{n=1}^{m}\tau
		\sum_{i=1}^{\infty}
		\norm{\bff g_{i,h}^{n-1}}{\bb L^2}^2
		\right]
		+
		C_p\bb E\left[
		\sum_{n=1}^{m}
		\left(
		\tau\sum_{i=1}^{\infty}
		\norm{\bff g_{i,h}^{n-1}}{\bb L^2}^2
		\right)^p
		\right].
	\end{align*}
	Since $\sum_{n=1}^{m}a_n^p
	\le
	\left(\sum_{n=1}^{m}a_n\right)^p$ for $a_n\ge0$,
	Young's inequality allows the first term on the right-hand side to be absorbed, yielding
	\begin{align}
		\label{equ:Q u p bound revised}
		\bb E[Q_{u,m}^p]
		\le
		C_p\bb E\left[
		\left(
		\sum_{n=1}^{m}\tau
		\sum_{i=1}^{\infty}
		\norm{\bff g_{i,h}^{n-1}}{\bb L^2}^2
		\right)^p
		\right]
		\le
		C_{p,T}
		\left[
		1+
		\sum_{n=1}^{m}\tau
		\bb E\left[(\mathcal E_h^{n-1})^p\right]
		\right],
	\end{align}
	where we also used \eqref{equ:g growth} and
	Lemma~\ref{lem:weighted-discrete-holder}.

	
	Repeating the same argument with
	$\norm{\curl\bff G_{B,h}^n}{\bb L^2}$ gives
	\begin{align}\label{equ:Q B p bound revised}
		\bb E[Q_{B,m}^p]
		&\le
		C_p\bb E\left[
		\left(
		\sum_{n=1}^{m}\tau
		\sum_{i=1}^{\infty}
		\norm{
			\curl\Pi_h^{\bb X}
			\bigl(
			\sigma\bff\chi_i
			+
			\eta\bff\chi_i\times\mathcal R_h\bff B_h^{n-1}
			\bigr)}
		{\bb L^2}^2
		\right)^p
		\right]
		\nonumber\\
		&\leq
		C_{p,T}
		\left[
		1+
		\sum_{n=1}^{m}\tau
		\bb E\left[(\mathcal E_h^{n-1})^p\right]
		\right]
		+
		C_p\delta_\chi^p
		\bb E\left[
		\left(
		\sum_{n=1}^{m}\tau
		\norm{\bff J_h^{n-1}}{\bb L^2}^2
		\right)^p
		\right],
	\end{align}
	where in the last step we used \eqref{equ:hall current coeff bound final}, Lemma~\ref{lem:weighted-discrete-holder}, and the fact that
	$\norm{\bff B_h^{n-1}}{\bb L^2}^2\le2\mathcal E_h^{n-1}$.
	The final term is retained until the conclusion of the proof, where it is absorbed into the dissipation moment.

	
	\emph{Step 6: Martingale estimate.}
	Let
	\[
	S_m:=\max_{0\le j\le m}\mathcal E_h^j.
	\]
	Since $\{M_h^n\}_{n=1}^{N}$ defined in
	\eqref{equ:mart increment revised} is a real-valued martingale difference
	sequence, we can apply Lemma~\ref{lem:discrete-BDG}. To this end, note that by \eqref{equ:mart increment revised} and the Cauchy--Schwarz inequality,
	\begin{align*}
		\abs{M_h^n}^2
		&\le
		2\norm{\bff G_{u,h}^n}{\bb L^2}^2
		\norm{\bff u_h^{n-1}}{\bb L^2}^2
		+
		2\norm{\curl\bff G_{B,h}^n}{\bb L^2}^2
		\norm{\bff B_h^{n-1}}{\bb L^2}^2
		\\
		&\le
		C\mathcal E_h^{n-1}
		\left(
		\norm{\bff G_{u,h}^n}{\bb L^2}^2
		+
		\norm{\curl\bff G_{B,h}^n}{\bb L^2}^2
		\right).
	\end{align*}
	Consequently,
	\[
	\sum_{n=1}^{m}\abs{M_h^n}^2
	\le
	CS_m\left(Q_{u,m}+Q_{B,m}\right).
	\]
	Using Lemma~\ref{lem:discrete-BDG}, Young's inequality, and
	\eqref{equ:Q u p bound revised}--\eqref{equ:Q B p bound revised}, we obtain
	\begin{align*}
		\bb E\left[
		\max_{1\le r\le m}
		\abs{\sum_{n=1}^{r}M_h^n}^{p}
		\right]
		&\le
		C_p\bb E\left[
		S_m^{p/2}
		\left(Q_{u,m}+Q_{B,m}\right)^{p/2}
		\right]
		\nonumber\\
		&\le
		\varepsilon\bb E[S_m^p]
		+
		C_{\varepsilon,p}
		\bb E\left[
		\left(Q_{u,m}+Q_{B,m}\right)^p
		\right]
		\nonumber\\
		&\le
		\varepsilon\bb E[S_m^p]
		+
		C_{\varepsilon,p,T}
		\left[
		1+
		\sum_{n=1}^{m}\tau
		\bb E\left[(\mathcal E_h^{n-1})^p\right]
		\right]
		+
		C_{\varepsilon,p}\delta_\chi^p
		\bb E\left[
		\left(
		\sum_{n=1}^{m}\tau
		\norm{\bff J_h^{n-1}}{\bb L^2}^2
		\right)^p
		\right].
	\end{align*}
	
	
	\emph{Step 7: Conclusion.}
	The case $p=1$ was established in Step~4, so let $p>1$. For
	$m=1,\ldots,N$, let $S_m$ be as in Step 6, and define
	\begin{align*}
		D_m &:=
		\sum_{n=1}^{m}
		\left(
		\norm{\bff u_h^n-\bff u_h^{n-1}}{\bb L^2}^2
		+
		\norm{\bff B_h^n-\bff B_h^{n-1}}{\bb L^2}^2
		\right)
		+
		\sum_{n=1}^{m}\tau
		\left(
		\nu\norm{\nabla\bff u_h^n}{\bb L^2}^2
		+
		\sigma\norm{\bff J_h^n}{\bb L^2}^2
		\right).
	\end{align*}
	Summing \eqref{equ:pathwise inequality revised} from $n=1$ to $m$,
	taking the maximum over the partial sums of the martingale term, raising
	to the power $p$, and using Steps~5--6 yields, for any sufficiently small
	$\varepsilon>0$,
	\begin{align}
		\label{equ:stability conclusion before absorption}
		\bb E[S_m^p]
		+
		c_p\bb E[D_m^p]
		&\le
		C_{p,T}
		\left(
		1
		+
		\bb E[(\mathcal E_h^0)^p]
		+
		\left(
		\sum_{n=1}^{N}
		\tau\norm{\bff f^n}{\bb H^{-1}}^2
		\right)^p
		\right)
		+
		C_{p,T}
		\sum_{n=1}^{m}
		\tau\bb E[(\mathcal E_h^{n-1})^p]
		+
		\varepsilon\bb E[S_m^p]
		\nonumber\\
		&\quad
		+
		C_{\varepsilon,p}\delta_\chi^p
		\bb E\left[
		\left(
		\sum_{n=1}^{m}
		\tau\norm{\bff J_h^{n-1}}{\bb L^2}^2
		\right)^p
		\right],
	\end{align}
	where $c_p>0$ is independent of $h$, $\tau$, and $m$.
	
	The shifted current term satisfies
	\begin{align*}
		\sum_{n=1}^{m}
		\tau\norm{\bff J_h^{n-1}}{\bb L^2}^2
		=
		\tau\norm{\bff J_h^0}{\bb L^2}^2
		+
		\sum_{n=1}^{m-1}
		\tau\norm{\bff J_h^n}{\bb L^2}^2
		\le
		\tau\norm{\bff J_h^0}{\bb L^2}^2
		+
		\frac{1}{\sigma}D_m.
	\end{align*}
	Consequently,
	\[
	\bb E\left[
	\left(
	\sum_{n=1}^{m}
	\tau\norm{\bff J_h^{n-1}}{\bb L^2}^2
	\right)^p
	\right]
	\le
	C_p\bb E\left[
	\left(
	\tau\norm{\bff J_h^0}{\bb L^2}^2
	\right)^p
	\right]
	+
	C_{p,\sigma}\bb E[D_m^p].
	\]
	We first choose $\varepsilon>0$ sufficiently small and then choose
	$\delta_p^\ast>0$ sufficiently small such that, whenever
	$\delta_\chi\le\delta_p^\ast$, the terms involving
	$\bb E[S_m^p]$ and $\bb E[D_m^p]$ on the right-hand side of
	\eqref{equ:stability conclusion before absorption} can be absorbed into
	the left-hand side. We thus obtain
	\begin{align*}
		\bb E[S_m^p]
		+
		\bb E[D_m^p]
		&\le
		C_{p,T}
		\left(
		1
		+
		\bb E[(\mathcal E_h^0)^p]
		+
		\bb E\left[
		\left(
		\tau\norm{\bff J_h^0}{\bb L^2}^2
		\right)^p
		\right]
		+
		\left(
		\sum_{n=1}^{N}
		\tau\norm{\bff f^n}{\bb H^{-1}}^2
		\right)^p
		\right)
		\\
		&\quad
		+
		C_{p,T}
		\sum_{n=1}^{m}
		\tau\bb E[(\mathcal E_h^{n-1})^p].
	\end{align*}
	Since $(\mathcal E_h^{n-1})^p\le S_{n-1}^p$, the discrete Gronwall
	lemma yields \eqref{equ:stoch stability revised main}. This completes the
	proof.
\end{proof}

\end{document}
